% !TeX root = ../../Book2.tex
\section{Hilbert 空间方法}
\label{b2:sec:1502}

上一章的极小化要求能量保留方程的全部信息。对非对称系数，这个要求可能失败；双线性型方法则直接研究等式 \(a(u,v)=\ell(v)\)。其中未知量占第一个位置，测试向量占第二个位置。本节把有界性和强制性分别转成算子的上界和下界，再用正交性证明可解性。

\subsection{双线性型}
\label{b2:subsec:150201}

设 \(H\) 是实 Hilbert 空间。\term[shuangxianxing-xing]{双线性型}[bilinear form]是对两个变量分别线性的映射 \(a:H\times H\to\mathbb R\)。对称性 \(a(u,v)=a(v,u)\) 是额外要求，本章的主要存在定理不作这一假设。例如
\[a(u,v)=\int_\Omega A\nabla u\cdot\nabla v\,\mathrm dx\]
在 \(H_0^1(\Omega)\) 上是双线性型；若矩阵 \(A(x)\) 对称，则型对称。一个纯反对称矩阵场对 \(a(u,u)\) 没有贡献，却可能改变 \(a(u,v)\)。

\ProseParagraphBreak

复 Hilbert 空间中通常改用对第一个变量线性、对第二个变量共轭线性的半双线性型。相应右端也对测试向量共轭线性，强制性写为 \(\operatorname{Re}a(u,u)\ge\alpha\lVert u\rVert^2\)。这与本书内积对第一变量线性的约定一致。为使证明中的对偶方向透明，下面先给出实空间版本，随后说明复版本的调整；不要把这一约定与分布配对的双线性约定混用。

\subsection{有界性}
\label{b2:subsec:150202}

双线性型的有界性是存在 \(M<\infty\)，使
\[|a(u,v)|\le M\lVert u\rVert_H\lVert v\rVert_H\qquad(u,v\in H).\]
它保证固定任一变量后所得泛函连续，也保证从稠密的测试子空间推广等式时极限合法。若 \(A\) 的矩阵算子范数几乎处处不超过 \(\Lambda\)，则
\[|a(u,v)|\le\Lambda\lVert\nabla u\rVert_2\lVert\nabla v\rVert_2.\]
这里有界性只使用系数的本性上界，不涉及其导数，也不要求点点连续。

\begin{proposition}[有界型的算子表示]
\label{b2:prop:15-form-operator}
设 \(a\) 是实 Hilbert 空间 \(H\) 上的有界双线性型。则存在唯一有界线性算子 \(A:H\to H\)，使
\[a(u,v)=\langle Au,v\rangle_H\qquad(u,v\in H),\]
且 \(\lVert A\rVert=\sup_{\lVert u\rVert,\lVert v\rVert\le1}|a(u,v)|\)。型对称当且仅当算子自伴。
\end{proposition}
\begin{proof}
固定 \(u\)，第一卷定理 8.4.7 把连续线性泛函 \(v\mapsto a(u,v)\) 表示为与唯一向量 \(Au\) 的内积。表示的唯一性和第一变量线性性给出 \(A\) 的线性性，且 \(\lVert Au\rVert\le M\lVert u\rVert\)。对单位向量分别取上确界给出范数等式。最后，\(a(u,v)=a(v,u)\) 等价于 \(\langle Au,v\rangle=\langle u,Av\rangle\)，正是自伴条件。
\end{proof}

这个 \(A\) 是依赖所选 Hilbert 内积的表示算子，不必等于微分表达式本身。例如在梯度内积的 \(H_0^1\) 中，Poisson 型 \(\int\nabla u\cdot\nabla v\) 的表示算子就是恒等算子；作为分布的 \(-\Delta\) 则把 \(H_0^1\) 映到 \(H^{-1}\)。分清值域有助于避免把有界算子与无界微分算子的谱性质混在一起。

\subsection{强制性}
\label{b2:subsec:150203}

称型 \(a\) 强制，是指存在 \(\alpha>0\)，使
\[a(u,u)\ge\alpha\lVert u\rVert_H^2\qquad(u\in H).\]
这里要求统一的 \(\alpha\)，而不只是每个非零向量上取正值。结合算子表示与 Cauchy--Schwarz 不等式，立刻得到
\[\alpha\lVert u\rVert_H\le\lVert Au\rVert_H.\]
它同时给出单射性和逆映射在值域上的有界性。不过，单有这种下界并不保证一般算子满射；例如 \(\ell^2\) 上的右移是等距单射却漏掉第一个坐标方向。Lax--Milgram 证明还要利用同一个型在测试向量处的正性，排除值域的正交补。

\ProseParagraphBreak

对椭圆型矩阵，条件 \(\xi^{\mathsf T}A(x)\xi\ge\lambda|\xi|^2\) 几乎处处成立时，梯度型在 \(H_0^1\) 的梯度范数下以 \(\alpha=\lambda\) 强制。在整个 \(H^1\) 上却并非如此，因为常函数在型中为零。选空间与选范数是定理假设的一部分，不能在写下“椭圆”之后省略核验。

\subsection{Lax–Milgram 定理}
\label{b2:subsec:150204}

\begin{theorem}[Lax--Milgram]
\label{b2:thm:15-lax-milgram}
设 \(H\) 是实 Hilbert 空间，\(a:H\times H\to\mathbb R\) 是双线性型，且对某些 \(M<\infty\)、\(\alpha>0\) 有
\[|a(u,v)|\le M\lVert u\rVert_H\lVert v\rVert_H,
 \qquad a(u,u)\ge\alpha\lVert u\rVert_H^2.\]
则对每个 \(\ell\in H^*\)，存在唯一 \(u\in H\) 满足
\[a(u,v)=\ell(v)\qquad(v\in H),\]
且 \(\lVert u\rVert_H\le\alpha^{-1}\lVert\ell\rVert_{H^*}\)。
\end{theorem}
\begin{proof}
用 \cref{b2:prop:15-form-operator} 表示 \(a(u,v)=\langle Au,v\rangle\)，并用 Riesz 表示 \(\ell(v)=\langle f,v\rangle\)。所求等式等价于 \(Au=f\)。强制性给出 \(\lVert Au\rVert\ge\alpha\lVert u\rVert\)，所以 \(A\) 单射。

若 \(Au_n\to y\)，下界说明 \((u_n)\) 是 Cauchy 序列，因而收敛到某个 \(u\in H\)。由 \(A\) 连续，\(Au=y\)，故值域闭。若 \(w\) 与值域正交，则 \(a(v,w)=\langle Av,w\rangle=0\) 对所有 \(v\in H\) 成立。取 \(v=w\)，强制性推出 \(w=0\)。因此值域的正交补为零；闭子空间的正交分解给出值域等于 \(H\)。于是 \(Au=f\) 唯一可解，且算子下界给出 \(\lVert u\rVert\le\alpha^{-1}\lVert f\rVert=\alpha^{-1}\lVert\ell\rVert\)。
\end{proof}

定理也称 \term[lax-milgram-dingli]{Lax--Milgram 定理}[Lax--Milgram theorem]。复空间版本把强制性改成实部下界，并把 \(\ell\) 视为连续共轭线性泛函；算子表示、闭值域和正交补的三步证明保持不变，在取 \(v=w\) 时使用 \(\operatorname{Re}a(w,w)\) 即可。

\ProseParagraphBreak

还可以把存在性看成一次收缩迭代。若 \(A\) 满足上述上下界，取 \(0<t<2\alpha/M^2\)，则
\[
 \lVert(I-tA)w\rVert^2
 =\lVert w\rVert^2-2t\langle Aw,w\rangle+t^2\lVert Aw\rVert^2
 \le(1-2t\alpha+t^2M^2)\lVert w\rVert^2.
\]
因此 \(u_{n+1}=u_n+t(f-Au_n)\) 收敛到唯一解。这个计算不仅提供另一种构造，也表明当 \(\alpha\) 很小而 \(M\) 很大时，允许的步长和收敛速度都会变差。定理的稳定性常数真实反映了问题接近失去强制性的程度。
